\chapter{Results} % Chapter 4
\label{Results}

% =============================================================================
% RESULTS (10-15 pages)
% Rubric: Results and Analysis (20% of Report)
% CRITICAL: Reports WHAT happened. Discussion chapter interprets WHY.
% =============================================================================

% -----------------------------------------------------------------------------
% SECTION 5.1: Experimental Setup Summary
% -----------------------------------------------------------------------------
\section{Experimental Overview}
% TODO: Brief recap of conditions actually run
%   - Total runs completed per phase
%   - Any deviations from planned design (budget constraints, model changes)
%   - Reference Methods chapter for full design details

% -----------------------------------------------------------------------------
% SECTION 5.2: OAT Parameter Characterization
% -----------------------------------------------------------------------------
\section{Parameter Characterization (OAT Sweeps)}

% TODO: 5.2.1 — Conflict parameters
\subsection{Conflict Parameters}
%   - conflict_cost_pct [2, 5, 10, 20]: effect on attack rate, Gini
%   - attack_take_pct [20, 40, 60, 80]: effect on attack rate, inequality
%   - Key finding: which parameters most strongly influence behavior?
%   - L1 vs L3 comparison: does reasoning depth modulate parameter sensitivity?

% TODO: 5.2.2 — Arming parameters
\subsection{Arming Parameters}
%   - arm_cost_pct [5, 10, 20]: effect on arm rate, defensive behavior
%   - arm_other_cost_pct [5, 10, 20]: effect on coalition formation
%   - Key finding: threshold effects? Phase transitions in arming behavior?

% TODO: 5.2.3 — Cooperation parameters
\subsection{Cooperation Parameters}
%   - invest_other_return_pct [10, 15, 25]: effect on cooperation rate
%     This is the most theoretically interesting: does increasing returns to
%     cooperation shift the system from conflict to cooperation?
%     Does reasoning depth modulate this shift?
%   - invest_other_cost_pct [5, 10, 20]: effect on cooperation cost sensitivity

% TODO: 5.2.4 — Spatial parameter
\subsection{Spatial Structure}
%   - interaction_radius [1, 2, 3]: effect on cooperation/conflict
%   - Expected: phase transition in cooperation at some radius threshold
%     (consistent with Perc, 2017; our Gemma 2 results)

% TODO: 5.2.5 — Toggle effects
\subsection{Toggle Effects}
%   - invest_self on/off: does removing self-investment force cooperation?
%   - memory on/off: does persistent memory change behavior vs stateless?
%     This is the first test of IV2 (information structure)

% TODO: 5.2.6 — OAT summary
\subsection{Summary: Sensitive Parameters}
%   - Table: parameter × effect size × reasoning depth interaction
%   - Identify top 2-3 parameters for factorial design
%   - Report partial η² for each parameter and reasoning level interaction

% -----------------------------------------------------------------------------
% SECTION 5.3: Reasoning Depth Effects (IV1)
% -----------------------------------------------------------------------------
\section{Reasoning Depth Effects}
% TODO: If separate reasoning depth sweep is run (all 4 levels × N reps)
%   - Action distributions per level (L0 vs L1 vs L2 vs L3)
%   - Gini trajectory per level
%   - Cooperation rate evolution per level
%   - Non-monotonic pattern: does L0→L1→L2→L3 show non-linear effects?
%   - Compare to Kuusela & Roy: do we see the Hobbesian Trap (more reasoning → more conflict)?
%   - Or is it more nuanced (L2 defensive, L3 exploitative)?
%
%   Statistical tests:
%   - One-way ANOVA: Gini ~ reasoning_level
%   - Pairwise comparisons with Bonferroni correction
%   - Effect sizes (Cohen's d) for each pair
%   - Bayes factors for key comparisons

% -----------------------------------------------------------------------------
% SECTION 5.4: Information Structure Effects (IV2)
% -----------------------------------------------------------------------------
\section{Information and Memory Effects}
% TODO: If memory sweep is run
%   - Memory ON vs OFF: behavioral differences
%   - Do agents with memory develop more stable relationships?
%   - Do agents with memory show more retaliation (they remember attacks)?
%   - Does memory reduce or increase inequality?
%
% TODO: Visibility effects
%   - Interaction of radius × memory: does memory matter more at small radius?
%   - hide_resources effect: does removing information about others change behavior?

% -----------------------------------------------------------------------------
% SECTION 5.5: Communication & Enforcement Effects (IV3)
% -----------------------------------------------------------------------------
\section{Communication and Enforcement Effects}
% TODO: If Phase 2d is completed
%   - No-comm vs cheap talk vs enforcement conditions
%   - Does cheap talk change behavior? (Crawford & Sobel: only if interests align)
%   - Does enforcement reduce conflict? (Hobbes's Leviathan hypothesis)
%   - Interaction with reasoning depth: does enforcement matter more at higher L?

% -----------------------------------------------------------------------------
% SECTION 5.6: Phase Transitions
% -----------------------------------------------------------------------------
\section{Phase Transitions}
% TODO: Phase transition analysis across IVs
%   - Gini vs control parameter (e.g., radius) per reasoning level
%   - Does the critical point shift with reasoning depth?
%   - EWS analysis: variance peak, autocorrelation increase before transition
%   - Network structure transitions: community count/stability vs parameter
%
% This is potentially the strongest result: if reasoning depth shifts the
% phase transition point, that demonstrates a genuine interaction between
% cognition and environmental structure.

% -----------------------------------------------------------------------------
% SECTION 5.7: Summary
% -----------------------------------------------------------------------------
\section{Summary of Results}
% TODO: Concise summary table mapping results to research questions
%   SQ1 (reasoning depth): [finding]
%   SQ2 (information): [finding]
%   SQ3 (communication): [finding]
%
% NOTE: Do NOT interpret findings here - that belongs in Discussion chapter.
